\section{Dimension-drop relations and Jiang--Su absorption}
\label{abs:section}

Proposition~\ref{prop:central-filler} supplies a matrix cone and an
element filling its defect. We first turn their relations into unital
maps from individual prime dimension-drop algebras. These maps need
not be compatible: nuclear approximation and bounded completely
positive lifts will instead produce a single unital map
$\Z\to F_\omega(A)$. Finally, half-flip unitaries and an approximate
intertwining give $A\cong A\otimes\Z$.

This route from central dimension-drop maps to absorption is
established. For separable, possibly nonunital algebras,
Nawata \cite[Proposition~5.1]{Nawata2013} states the criterion using
the sequence quotient. It fits the strongly self-absorbing framework
of Toms--Winter \cite{TW2007}; a formulation for arbitrary strongly
self-absorbing algebras is given in
\cite[Theorem~2.2]{PTV2025}. We give the lifting and intertwining
steps for the fixed ultrafilter quotient used here.

\subsection{The unit in the annihilator quotient}

Put
\[
 \Ann(A)=\{x\in A_\omega:xa=ax=0\text{ for every }a\in A\},
 \qquad
 F_\omega(A)=(A_\omega\cap A')/\Ann(A).
\]
The annihilator is a closed two-sided ideal of $A_\omega\cap A'$.
It is contained in that relative commutant, since all its products
with $A$ vanish.
If $(e_n)$ is a positive contractive sequential approximate unit of
$A$, then $e=[(e_n)]\in A_\omega\cap A'$, and its class is the unit
of $F_\omega(A)$.  Indeed $ea=ae=a$ for $a\in A$, while for
$x\in A_\omega\cap A'$ the differences $ex-x$ and $xe-x$
annihilate $A$ on both sides.  For example
$(ex-x)a=(ea-a)x=0$, and the other products follow in the same way.

Apply Proposition~\ref{prop:central-filler}.  The selfadjoint element
$S^*S+h-e$ annihilates $A$ on the right by
\eqref{cen:filler-relations}, and therefore on the left by taking
adjoints.  Writing $s$ for the quotient image of $S$, and using the
same letters $\alpha,h$ for the quotient images of the map and its
weight, gives the exact relations
\begin{equation}\label{abs:quotient-relations}
 s^*s=1-h,\qquad \alpha(e_{11})s=s,
 \qquad h=\alpha(1).
\end{equation}
In particular $s$ is a contraction in the quotient.  The argument uses
the quotient unit represented by $(e_n)$, without requiring a unit in
$A$.

\subsection{An explicit dimension-drop model}

For positive integers $r,s$, write
\[
 I(r,s)=\{f\in C([0,1],M_r\otimes M_s):
 f(0)\in M_r\otimes1_s,\ f(1)\in1_r\otimes M_s\}.
\]

Relations~\eqref{abs:quotient-relations} are the order-zero form of the dimension-drop
presentation in \cite[Proposition~2.1]{Sato2010}, following
\cite[Proposition~5.1, (iii)$\Rightarrow$(iv)]{RW2010}. We give an explicit model and verify
that every representation of the relations factors through it.

\begin{proposition}\label{prop:dimensiondrop}
Let $B$ be a unital $C^*$-algebra, let $p\ge2$, and suppose that a
\cpc order-zero map $\alpha:M_p\to B$ and $s\in B$ satisfy
\[
 h=\alpha(1),\qquad s^*s=1-h,\qquad \alpha(e_{11})s=s.
\]
There is a unital homomorphism
$I(p,p+1)\to B$ whose range contains $\alpha(M_p)$ and $s$.
Consequently $F_\omega(A)$ receives a unital homomorphism from
$I(p,p+1)$ for every $p\ge2$.
\end{proposition}

\begin{proof}
\emph{The model and its generators.}
Decompose the first
matrix leg of $M_{p+1}\otimes M_p$ into blocks of sizes $p$ and $1$,
and put
\begin{align}
 A_a(t)&=\begin{pmatrix}a\otimes1_p&0\\0&ta\end{pmatrix},
       &&a\in M_p,\label{abs:model-map}\\
 T(t)&=\sqrt{1-t}\ e_{1,p+1}^{(p+1)}\otimes1_p.
       \label{abs:model-element}
\end{align}
The upper diagonal block has size $p^2$ and the lower one size $p$.
At each $t$ the map $a\mapsto A_a(t)$ is completely positive,
contractive, and order zero.  Direct matrix multiplication gives
\begin{equation}\label{abs:model-relations}
 T^*T=1-A_1,\qquad A_{e_{11}}T=T,
\end{equation}
and
\begin{equation}\label{abs:scalar-identity}
 T^*T+\sum_{i=1}^p A_{e_{i1}}TT^*A_{e_{1i}}
           =(1-t)1_{p(p+1)}.
\end{equation}
For the last identity, $T^*T$ is $(1-t)$ times the bottom block
projection, and the $i$th summand is $(1-t)$ times the $i$th top
block projection.

Let
\[
 j(a)=\operatorname{diag}(a\otimes1_p,a)
\]
and define the boundary algebra
\[
 B_p=\{f\in C([0,1],M_{p+1}\otimes M_p):
 f(0)\in M_{p+1}\otimes1_p,\ f(1)\in j(M_p)\}.
\]
We claim that $A_a$ and $T$ generate $B_p$.
First $A_1+T^*T=1$, and \eqref{abs:scalar-identity} gives the
coordinate scalar function $t1$.  Thus their generated algebra
contains $C([0,1])1$.
At $t=1$ the generators have image $j(M_p)$.
At $0\le t<1$, the nonzero link $T(t)$, its adjoint, and the top
matrix action yield all matrix units of $M_{p+1}\otimes1_p$:
the links $A_{e_{i1}}(t)T(t)$ give the $i$th top-to-bottom units,
and their products give all top-to-top units and the bottom
projection.  At $t=0$ these are exactly the prescribed fiber.
When $0<t<1$, compression of $A_a(t)$ to the bottom block and
division by the scalar $t$ gives the full bottom copy of $M_p$.
Multiplying by the preceding matrix units gives all of
$M_{p+1}\otimes M_p$ in that fiber.

Fiber generation and the scalar functions prove the claim in norm.
More explicitly, for $f\in B_p$ and $\eps>0$, at each $t$ choose
an element $g_t$ of the generated algebra whose value at $t$ agrees
with $f(t)$, or is within $\eps/2$ of it.  Continuity gives an open
neighborhood on which $\|g_t-f\|<\eps$.  A finite subcover and a
subordinate scalar partition of unity produce an element of the
generated algebra within $\eps$ of $f$ on the whole interval.

The representation $j$ of $M_p$ has multiplicity $p+1$, so some
unitary $V\in M_{p+1}\otimes M_p$ conjugates $j(M_p)$ onto
$1_{p+1}\otimes M_p$.  The finite-dimensional unitary group is
path connected; choose a unitary path $V(t)$ from $1$ to $V$.
Conjugation by this path changes only the endpoint presentation:
the boundary algebra becomes the one with fibers
$M_{p+1}\otimes1_p$ at $0$ and $1_{p+1}\otimes M_p$ at $1$.
Reversing the interval and interchanging the tensor legs identifies
it with $I(p,p+1)$.

\emph{Arbitrary representations of the relations.}
We have identified the algebra generated by the model. To obtain the
asserted map into $B$, we now show that every pair $(\alpha,s)$ in the
statement gives a representation of that algebra. Represent $B$ faithfully and
unitally on a Hilbert space $H$.  The order-zero structure theorem
gives a supporting representation $\rho:M_p\to\mathcal B(H)$,
with unit the support projection of $h=\alpha(1)$, such that
$\alpha(a)=h\rho(a)=\rho(a)h$.
In the polar decomposition
\[
 s=u(1-h)^{1/2},
\]
the initial space of $u$ is
$H_s=1_{[0,1)}(h)H$.  The equation $\alpha(e_{11})s=s$ implies
that $\alpha(e_{11})$ is the identity on the closed range of $s$.
Since $0\le\alpha(e_{11})\le h\le1$, both $h$ and
$\rho(e_{11})$ are the identity on this range.  Thus $u$ is an
isometry from $H_s$ into
$1_{\{1\}}(h)\rho(e_{11})H$.

For $1\le i\le p$ the subspaces
$\rho(e_{i1})uH_s$ are mutually orthogonal.  Their direct sum
$H_t$ reduces $\rho$ and is identified with $\C^p\otimes H_s$
by the unitary sending $e_i\otimes\zeta$ to
$\rho(e_{i1})u\zeta$.  On $H_t$ the weight is $1$ and the
representation is $a\otimes1_{H_s}$.  Since $H_t$ is contained
in the $h=1$ space, it is orthogonal to $H_s$.
The orthogonal complement
\[
 H_r=1_{\{1\}}(h)H\ominus H_t
\]
reduces both generators: $s$ and $s^*$ vanish there, and $\alpha$
acts there as the unital matrix representation $\rho|_{H_r}$.
This summand is therefore obtained by evaluating $B_p$ at $1$.

On $H_t\oplus H_s\cong\C^{p+1}\otimes H_s$ the generators
are precisely
\begin{equation}\label{abs:arbitrary-representation}
 \alpha(a)=
 \begin{pmatrix}a\otimes1_{H_s}&0\\0&h_s\rho_s(a)\end{pmatrix},
 \qquad
 s=e_{1,p+1}^{(p+1)}\otimes(1-h_s)^{1/2},
\end{equation}
where $h_s$ and $\rho_s$ denote the source restrictions, with
$\rho_s=0$ on $\ker h_s$.
Split $H_s$ as
$H_0\oplus H_{\mathrm{mid}}$ according to the spectral sets
$\{0\}$ and $(0,1)$.  On $\C^{p+1}\otimes H_0$,
\eqref{abs:arbitrary-representation} is the amplification of
evaluation at $0$: its bottom matrix action is zero and its link
has coefficient $1$.
On $H_{\mathrm{mid}}$ the representation $\rho_s$ is unital and
commutes with $h_s$.  Continuous functional calculus therefore gives
a unital representation
\[
 C([0,1])\otimes M_p\longrightarrow\mathcal B(H_{\mathrm{mid}}),
 \qquad f\otimes a\longmapsto f(h_s)\rho_s(a).
\]
Amplifying by $M_{p+1}$ and restricting to $B_p$ sends
\eqref{abs:model-map}--\eqref{abs:model-element} to
\eqref{abs:arbitrary-representation} on this summand.
Together with the two endpoint representations this gives a unital
representation of $B_p$ on all of $H$, sending its generators to
$\alpha(a)$ and $s$.  Since these generators generate $B_p$, the
image lies in the represented copy of $B$.  The identification of
$B_p$ with $I(p,p+1)$ proves the assertion.
\end{proof}

\subsection{From prime blocks to a central copy of \texorpdfstring{$\Z$}{Z}}

We use the standard structural properties of the Jiang--Su algebra.
It is separable, unital, and nuclear, and is an injective unital
inductive limit of prime dimension-drop blocks
\cite[Proposition~2.5 and its proof]{JS1999}.
It is also strongly self-absorbing
\cite[Theorems~7.6 and~8.7]{JS1999}; see the framework in
\cite{TW2007} and the later construction in \cite{RW2010}.
We next give the passage from the
individual block maps to a homomorphism of the limit, since the block
maps obtained above need not be compatible.

\begin{lemma}\label{abs:prime-block-map}
If $r$ and $s$ are relatively prime, there is a unital homomorphism
$I(r,s)\to I(p,p+1)$ for some $p\ge2$.
\end{lemma}

\begin{proof}
The congruences $p\equiv0\pmod r$ and $p\equiv-1\pmod s$
have a simultaneous solution, and adding a sufficiently large multiple
of $rs$ makes it at least $2$.  Write $p=mr$, $p+1=ns$, and use
the unital amplifications
\[
 M_r\longrightarrow M_p,\quad a\longmapsto a\otimes1_m,
 \qquad
 M_s\longrightarrow M_{p+1},\quad b\longmapsto b\otimes1_n.
\]
Their tensor product, applied pointwise, sends the $0$ boundary
$M_r\otimes1_s$ into $M_p\otimes1_{p+1}$ and the $1$ boundary
$1_r\otimes M_s$ into $1_p\otimes M_{p+1}$.
\end{proof}

The next lemma lets us approximately retract a finite subset of $\Z$
into a sufficiently late stage. Composing the retraction with any
homomorphism from that stage gives a completely positive map defined
on all of $\Z$. Approximation on the prescribed elements and their
products makes this map approximately multiplicative.

\begin{lemma}\label{abs:approximate-retraction}
Let $D\subset L$ be a unital inclusion with $D$ nuclear.  Given a
finite subset $G\subset D$ and $\eps>0$, there is a \cpc map
$R:L\to D$ with $\|R(d)-d\|<\eps$ for $d\in G$.
\end{lemma}

\begin{proof}
Choose a nuclear approximation
$D\xrightarrow{\theta}M_k\xrightarrow{\eta}D$, with both maps
completely positive and contractive, whose composition is within
$\eps$ of the identity on $G$.  The matrix-valued extension theorem
of Arveson \cite[Theorem~1.2.3]{Arveson1969} extends $\theta$ to a \cpc map
$\widetilde\theta:L\to M_k$.  Then
$R=\eta\widetilde\theta$ has the stated properties.  Applying the
approximation also to $1_D$ makes $R(1_L)$ as close to $1_D$ as
desired; an exactly unital retraction is unnecessary.
\end{proof}

\begin{proposition}\label{prop:Z-map}
Let $A$ be a nonzero separable nuclear $C^*$-algebra satisfying
\textup{(C)}, and fix any free ultrafilter $\omega$ on $\mathbb N$.
Then there is a unital homomorphism $\Z\to F_\omega(A)$.
\end{proposition}

\begin{proof}
Realize $\Z$ as the closure of an increasing sequence of its
injectively included prime dimension-drop stages $D_l$.
By Proposition~\ref{prop:dimensiondrop} and
Lemma~\ref{abs:prime-block-map}, every $D_l$ has a unital
homomorphism $\chi_l:D_l\to F_\omega(A)$.

Let $\mathcal Z_0$ be a countable dense $*$-subalgebra over
$\mathbb Q+i\mathbb Q$ containing $1$, and exhaust it by finite
sets.  For any finite set of contractions $F\subset\mathcal Z_0$
and any $\eps>0$, approximate $F\cup F^2\cup\{1\}$ by one
stage, and apply Lemma~\ref{abs:approximate-retraction} to those
stage approximants.  With the errors chosen sufficiently small,
the resulting \cpc map $R:\Z\to D_l$, viewed in $\Z$, satisfies
\[
 \|R(x)-x\|<\eps
 \quad(x\in F\cup F^2\cup\{1\}).
\]
Consequently, for $x,y\in F$,
\[
 \|R(xy)-R(x)R(y)\|
 \le\|R(xy)-xy\|+\|xy-R(x)R(y)\|<3\eps.
\]
Here the norm calculation takes place in $\Z$, with $D_l$ included
isometrically.  Composing with $\chi_l$ produces \cpc maps
$\beta_n:\Z\to F_\omega(A)$ such that
\begin{equation}\label{abs:quotient-approximation}
 \|\beta_n(xy)-\beta_n(x)\beta_n(y)\|\longrightarrow0,
 \qquad \|\beta_n(1)-1\|\longrightarrow0
\end{equation}
for every $x,y\in\mathcal Z_0$.  Including successively larger
scaled tests handles elements outside the unit ball.

We now turn these approximate quotient relations into exact ones.
Put $C_A=A_\omega\cap A'$, and let $q:C_A\to F_\omega(A)$
be the quotient map.  By the nuclear completely positive lifting
theorem of Choi--Effros
\cite[Theorem~3.10 and Corollary~3.11]{CE1976}, each $\beta_n$ has a \cpc
lift $L_n:\Z\to C_A$.  Regarding $L_n$ as a map into $A_\omega$
and applying the same theorem to
$\ell^\infty(\N,A)\to A_\omega$ gives a \cpc lift
\[
 \widetilde L_n:\Z\longrightarrow\ell^\infty(\N,A),
 \qquad \widetilde L_n(x)=(L_{n,k}(x))_{k\ge1}.
\]
Each coordinate map $L_{n,k}:\Z\to A$ is completely positive
and contractive, so all its values on each fixed test have a bound
independent of $n$ and $k$.

The relevant quotient norm errors are multiplied errors.
If $d\in C_A$ and $a\in A$, then
\begin{equation}\label{abs:multiplied-quotient-norm}
 \|da\|,\ \|ad\|\le\|q(d)\|\,\|a\|.
\end{equation}
Indeed, replacing $d$ by $d+j$ with $j\in\Ann(A)$ does not
change either product, and taking the infimum over such $j$ gives
the inequalities.  Apply this to
\[
 d=L_n(xy)-L_n(x)L_n(y),\qquad d=L_n(1)-e,
\]
where $e$ represents the unit as above.  Equations
\eqref{abs:quotient-approximation} and
\eqref{abs:multiplied-quotient-norm} give errors tending to zero
in $A_\omega$ after multiplication by each fixed $a\in A$.
Also $[L_n(x),a]=0$ exactly for every $n,x,a$.

Choose a countable dense set $\{a_j\}$ in $A$, and diagonalize
over $\mathcal Z_0$, its product tests, and these $a_j$.
More explicitly, choose increasing row indices $m_l$ so that the
first $l$ multiplied defects in row $m_l$ have ultrapower norm less
than $2^{-l-1}$. Their coordinate errors, together with the first
$l$ commutator errors, are less than $2^{-l}$ on some set
$E_l\in\omega$; replacing these sets by their successive finite
intersections with $\{k\ge l\}$ makes them decreasing. Let $l(k)$
be the largest $l\le k$ with $k\in E_l$, or zero when no such level
exists. Put $\gamma_k=L_{m_{l(k)},k}$ if $l(k)>0$, and $\gamma_k=0$
otherwise. Then $l(k)\to_\omega\infty$, and
\[
 \Gamma(x)=[(\gamma_k(x))_k]
\]
defines a \cpc map $\Gamma:\Z\to A_\omega\cap A'$. The maps
$\gamma_k$ are linear and completely positive at every coordinate; no
separate lifting of individual elements is involved.  The diagonal
conditions say
\[
 (\Gamma(xy)-\Gamma(x)\Gamma(y))a=0,\qquad
 (\Gamma(1)-e)a=0
 \quad(x,y\in\mathcal Z_0,\ a\in A).
\]
All the displayed defects are central relative to $A$, so they
annihilate it on both sides.  Passing through $q$ therefore gives
multiplicativity and unitality on $\mathcal Z_0$.  Continuity extends
these identities to $\Z$ and proves the proposition.
\end{proof}

\subsection{Connected half-flip unitaries and their lifts}

\begin{lemma}\label{abs:connected-unitaries}
The unitary groups of $\Z$ and $\Z\otimes\Z$ are path connected.
If $q:D\to D/J$ is a unital quotient map, every unitary in the
identity component of $\mathcal U(D/J)$ lifts to a unitary in $D$.
\end{lemma}

\begin{proof}
We first check connectedness for a prime block $I(r,s)$, with
$\gcd(r,s)=1$.  A unitary is a continuous path in $U(rs)$ whose
endpoints belong to
\[
 H_0=U(r)\otimes1_s,\qquad H_1=1_r\otimes U(s).
\]
Both endpoint groups are connected.  Attach a path in $H_0$ from
$1$ to the initial value and a path in $H_1$ from the final value
to $1$.  Attaching these paths gradually, while reparametrizing the
original path, is a homotopy through paths satisfying the same
endpoint conditions.  The result is a based loop in $U(rs)$.
Changing the attached $H_0$ path by an endpoint loop changes its
determinant winding by an arbitrary multiple of $s$; changing the
$H_1$ path changes it by an arbitrary multiple of $r$.  Bezout's
identity therefore permits choices for which the total winding is
zero.  The determinant identifies $\pi_1(U(rs))$ with $\mathbb Z$,
so this loop contracts with its basepoint fixed.  That contraction
retains endpoint value $1$, and hence remains in the prime block.
Thus every unitary of $I(r,s)$ is connected to $1$.

For a unitary in the injective unital inductive limit $\Z$, choose
a sufficiently close element of a stage.  Its two unitary defects
are small, so its polar part is a unitary in the same stage and is
still close to the original unitary.  Two unitaries of distance less
than $1$ are joined by the logarithm path for their quotient.
Connectedness of the stage unitary group proves connectedness of
$\mathcal U(\Z)$.  Strong self-absorption gives an isomorphism
$\Z\otimes\Z\cong\Z$, proving the same assertion for the tensor
square.

Finally, a unitary $v$ connected to $1$ in $D/J$ is a finite product
of exponentials of selfadjoints.  To see this, subdivide a continuous
unitary path finely enough that consecutive values differ by less
than $1$, and take a selfadjoint logarithm of each consecutive
quotient.  Lift these finitely many selfadjoints through $q$ and
multiply their exponentials in $D$.  The resulting unitary lifts $v$.
\end{proof}

\begin{lemma}\label{abs:finite-unitary-criterion}
Let $A$ be a separable $C^*$-algebra, and suppose there is a unital
homomorphism $\varphi:\Z\to F_\omega(A)$.
For finite sets $F\subset A$, $G\subset A\otimes\Z$, and
$\eps>0$, there is a unitary $u\in A^{\sim}\otimes\Z$ such that
\begin{equation}\label{abs:unitary-criterion}
 \|[u,a\otimes1]\|<\eps\quad(a\in F),\qquad
 \dist(u^*bu,A\otimes1)<\eps\quad(b\in G).
\end{equation}
\end{lemma}

\begin{proof}
Write $D=A^{\sim}\otimes\Z$ and $A_1=A\otimes1$, and form the
unital algebra and its annihilator quotient
\[
 C_D=D_\omega\cap A_1',\qquad
 Q=C_D/\Ann_{D_\omega}(A_1).
\]
First-factor sequences give a homomorphism
$F_\omega(A)\to Q$: centrality and annihilation are both preserved
by tensoring with $1$.  This homomorphism is unital, since the
approximate-unit representative $e\otimes1$ differs from $1_{D_\omega}$
by an annihilator of $A_1$.  Its composition with $\varphi$ is a
unital map $\varphi_1:\Z\to Q$ and commutes with the constant
second-factor map $\varphi_2(z)=1\otimes z+\Ann(A_1)$.
Because $\Z$ is nuclear, these two maps give a unital homomorphism
\[
 \Theta:\Z\otimes\Z\longrightarrow Q,
 \qquad \Theta(z\otimes z')=\varphi_1(z)\varphi_2(z').
\]

The approximately inner half-flip of the strongly self-absorbing
algebra $\Z$ \cite[Proposition~1.5]{TW2007} says that, for a finite set $H\subset\Z$
and $\delta>0$, there is a unitary $V\in\Z\otimes\Z$ with
\begin{equation}\label{abs:half-flip}
 \|V^*(1\otimes z)V-z\otimes1\|<\delta\quad(z\in H).
\end{equation}
By Lemma~\ref{abs:connected-unitaries}, $V$ is connected to $1$.
Thus $\Theta(V)$ is connected to $1$ in $Q$, and the same lemma
lifts it to a unitary $W\in C_D$.

This unitary has a representing sequence of actual unitaries in $D$.
Indeed choose bounded representatives $x_k$.  Since $W$ is unitary,
both $\|x_k^*x_k-1\|$ and $\|x_kx_k^*-1\|$ tend to zero along
$\omega$.  On a good set both are less than $1/2$, and the polar
parts $u_k=x_k(x_k^*x_k)^{-1/2}$ are unitaries with
$\|u_k-x_k\|\to_\omega0$.  Put $u_k=1$ on the remaining
coordinates.  Then $W=[(u_k)]$, and $W\in C_D$ gives
\begin{equation}\label{abs:lift-centrality}
 \|[u_k,a\otimes1]\|\longrightarrow_\omega0\quad(a\in A).
\end{equation}

For each $z\in H$ choose a lift $c_z\in A_\omega\cap A'$ of
$\varphi(z)$; a single \cpc lift of $\varphi$ is available by
Choi--Effros lifting, so these lifts may satisfy
$\|c_z\|\le\|z\|$.  In $C_D$ the difference
\[
 d_z=W^*(1\otimes z)W-c_z\otimes1
\]
has quotient norm less than $\delta$ by \eqref{abs:half-flip}.
The analogue of \eqref{abs:multiplied-quotient-norm} gives
\[
 \|(a\otimes1)d_z\|\le\delta\|a\|.
\]
Since $W$ commutes with $a\otimes1$ in the ultrapower, this is
the concrete norm estimate
\begin{equation}\label{abs:half-flip-multiplied}
 \|W^*(a\otimes z)W-(ac_z)\otimes1\|\le\delta\|a\|.
\end{equation}
Here $(ac_z)\otimes1$ is represented by a first-factor sequence
$(ac_{z,k})\otimes1$, with $ac_{z,k}\in A$.
Thus no estimate in an annihilator quotient is being used as an
unmultiplied estimate in a norm ultrapower.

For finitely many elementary tensors, choose one coordinate where
\eqref{abs:lift-centrality} and
\eqref{abs:half-flip-multiplied} all have the desired accuracy.
This gives \eqref{abs:unitary-criterion} for those tensors.
For a general finite $G$, approximate each of its elements by a
finite sum of elementary tensors, include all the corresponding
second factors in $H$, and choose $\delta$ small enough for the
finite sums of the bounds $\delta\|a\|$.  Unitary conjugation
preserves the initial approximation errors, so the same conclusion
holds for $G$.
\end{proof}

\subsection{The final intertwining}

The construction of the central map used nuclearity and condition
\textup{(C)}. The implication from that map to absorption requires
only separability, as the following proposition makes explicit.

\begin{proposition}\label{prop:absorption}
If a separable $C^*$-algebra $A$ admits a unital homomorphism
$\Z\to F_\omega(A)$, then $A\cong A\otimes\Z$.

\end{proposition}

\begin{proof}
The zero algebra is immediate, so assume $A\ne0$.
Set $B=A\otimes\Z$ and let $\iota(a)=a\otimes1$.
The algebra $A^{\sim}\otimes\Z$ acts by multipliers on its ideal
$B$, so all the unitaries in Lemma~\ref{abs:finite-unitary-criterion}
act as automorphisms of $B$ by conjugation.

Choose sequences $(a_j)$ and $(b_j)$ dense in $A$ and $B$,
respectively, and put $\eps_n=2^{-n}$.
We construct products of multiplier unitaries $U_n$ and embeddings
$\iota_n=\Ad(U_n)\circ\iota$ as follows.  Start with $U_0=1$.
Before step $n$, let $F_n$ be a finite subset of $A$ containing
$a_1,\ldots,a_n$ and all elements $c_{k,j}$ chosen at earlier
steps $k<n$, $1\le j\le k$.  Apply
Lemma~\ref{abs:finite-unitary-criterion} with these domain tests and
the finite target set
\[
 \{U_{n-1}^*b_jU_{n-1}:1\le j\le n\}.
\]
Obtain $u_n$ with errors less than $\eps_n$, and put
$U_n=U_{n-1}u_n$.
The commutator estimate gives
\begin{equation}\label{abs:successive-embeddings}
 \|\iota_n(a)-\iota_{n-1}(a)\|<\eps_n
 \quad(a\in F_n),
\end{equation}
and the range estimate gives elements $c_{n,j}\in A$ with
\begin{equation}\label{abs:stored-preimages}
 \|\iota_n(c_{n,j})-b_j\|<\eps_n
 \quad(1\le j\le n).
\end{equation}
Store these finitely many preimages in every subsequent $F_m$.
They are legitimate finite norm tests: each has finite norm, indeed
$\|c_{n,j}\|<\|b_j\|+\eps_n$, because $\iota_n$ is isometric.

For each dense domain element $a_j$, equation
\eqref{abs:successive-embeddings} holds at every step $n\ge j$.
Summability of $(\eps_n)$ gives convergence of $\iota_n(a_j)$.
As all $\iota_n$ are isometric, approximation by the dense set
extends this to point-norm convergence on $A$.
The limit $\Phi:A\to B$ is a homomorphism by continuity of
multiplication and adjoints, and it is isometric since
$\|\Phi(a)\|=\lim_n\|\iota_n(a)\|=\|a\|$.

Fix $j$ and take any $n\ge j$.  The stored preimage $c_{n,j}$
belongs to every $F_m$ with $m>n$, so
\[
 \|\Phi(c_{n,j})-b_j\|
 \le\eps_n+\sum_{m>n}\eps_m=2^{1-n}.
\]
Letting $n\to\infty$ shows that every $b_j$ is in the closure
of $\Phi(A)$.  The range of an isometric homomorphism is closed,
and the $b_j$ are dense, so $\Phi(A)=B$.  This proves the desired
isomorphism.


\end{proof}
